\chapter{Open Gaps and Future Work}
\label{sec:gaps}
% ===================================================================

This skeleton leaves several arguments at the level of conjecture
or intuition.
We list the main gaps here, both as an honest accounting and as a
research program.

\begin{enumerate}[label=\textbf{Gap \arabic*.}, leftmargin=*, itemsep=8pt]

  \item \textbf{Communication degradation (Conjecture~\ref{conj:degradation}).}
        The claim that no $\epsilon$-coherent path can span a
        population of quark scale requires a precise model of the
        communication geometry of a generic interactive GSLT, and
        a proof that diameter grows faster than the coherence
        length.
        The argument is information-theoretic and should be
        provable, but the details depend on the amplitude structure
        of the specific GSLT.

  \item \textbf{Error-correction overhead.}
        The claim that error-correction cannot bootstrap coherence
        at sufficient scale is currently informal.
        A precise argument would model error-correction as a
        sub-computation in the GSLT, account for its resource
        consumption, and show that the overhead exceeds available
        bandwidth beyond a critical scale.
        This is analogous to the no-cloning theorem in quantum
        information, and a similar proof strategy may apply.

  \item \textbf{Existence of prime combinator sets (Conjecture).}
        It remains to show that every interactive GSLT with a
        finitely generated grammar admits a prime combinator set.
        The proof strategy would likely use a well-founded
        induction on interaction complexity.

  \item \textbf{Relationship between metric and topological clusters.}
        The communication-radius clusters (metric) and the
        coherence clusters (topological/compact) are defined
        differently.
        A theorem identifying conditions under which they coincide
        would unify the two halves of the clustering argument.

  \item \textbf{Convergence of the world model (Proposition~\ref{prop:false-ontology}).}
        The claim that the agent's world model converges to the
        nerve fragment visible from its cluster requires a precise
        convergence theorem for the scientific method iteration
        under the combined constraints of finite budget, bounded
        coherence radius, and cluster compactness.

  \item \textbf{The recursive hierarchy.}
        The ontological hierarchy is defined, but it is not shown
        that the hierarchy is well-founded (terminates at a fixed
        point) or infinite.
        In physics, the hierarchy appears to terminate at the
        Planck scale; the computational analogue of a Planck-scale
        termination condition would be a significant result.

  \item \textbf{Stay's functorial construction.}
        The argument that agents can manipulate
        HML formulae as first-class terms relies on Stay's
        functorial generalization of the lambda cube to GSLTs.
        A full treatment of this construction, and in particular
        the canonical embedding back into the original GSLT and
        its interaction with the cluster topology, is deferred
        to a subsequent paper.

\end{enumerate}

% ===================================================================
